% !TEX root = ../../main.tex

\chapter{Extended OPE lines and the shadow comparison problem}
\label{chap:shadow-tower-extended-families}

\section*{Line-restricted data}
\label{sec:shadow-tower-extended-overview}

For a strongly generated vertex algebra, an OPE line records singular
products involving one selected generator.  A shadow coefficient is a
different object: it is obtained from the full ordered Maurer--Cartan
tensor by a specified projection after the configuration-space residue
maps have been constructed.  This chapter keeps these two levels
separate for $\mathcal W_3$, Bershadsky--Polyakov,
$\mathcal W_\infty$, and a super-Yangian test line.

The common type signature is
\[
 (\text{Open quadrant},\ \text{OPE line presentation},\
 \text{Beilinson levels }1\text{--}3,\
 H_{\mathrm{line}}+H_{\mathrm{bar}}+H_{\mathrm{proj}}).
\]
Here $H_{\mathrm{line}}$ fixes the two-point form and generator
normalization, $H_{\mathrm{bar}}$ constructs the signed residue
differential, and $H_{\mathrm{proj}}$ supplies a chain map from the
multi-generator tensor to the scalar line.

\section{Zamolodchikov norms and the $\mathcal W_3$ OPE}
\label{sec:zamolodchikov-norms}

Set
\[
 \Lambda={:TT:}-\frac{3}{10}\partial^2T.
\]

\begin{proposition}[Canonical two-point norms;
\ClaimStatusProvedElsewhere]
\label{prop:canonical-two-point-norms}
In Zamolodchikov normalization,
\begin{equation}\label{eq:zamolodchikov-norms}
 N_T=\langle T,T\rangle=\frac c2,
 \qquad
 N_W=\langle W,W\rangle=\frac c3,
 \qquad
 N_\Lambda=\langle\Lambda,\Lambda\rangle
 =\frac{c(5c+22)}{10}.
\end{equation}
\end{proposition}

\begin{proof}
The first two identities are the leading coefficients in the $TT$ and
$WW$ OPEs.  For the third, the level-four vacuum Gram matrix in the
basis $(L_{-2}^2\mathbf1,L_{-4}\mathbf1)$ is
\[
 \begin{pmatrix}
  c(c+8)/2 & 3c\\
  3c & 5c
 \end{pmatrix}.
\]
The vector of $\Lambda$ is $(1,-3/5)$, and direct multiplication gives
$c(5c+22)/10$; see Zamolodchikov~\cite{Zam85} and the review
\cite{BowSch93}.
\end{proof}

\begin{proposition}[Zamolodchikov-normalized $WW$ OPE;
\ClaimStatusProvedElsewhere]
\label{prop:fateev-lukyanov-alpha}
The singular part of the $\mathcal W_3$ OPE is
\[
\begin{aligned}
W(z)W(w)\sim{}&\frac{c/3}{(z-w)^6}
 +\frac{2T}{(z-w)^4}+\frac{\partial T}{(z-w)^3}\\
&+\frac{\frac{3}{10}\partial^2T+\alpha\Lambda}{(z-w)^2}
 +\frac{\frac{1}{15}\partial^3T+\frac{\alpha}{2}\partial\Lambda}{z-w},
\end{aligned}
\]
where
\begin{equation}\label{eq:alpha-coefficient}
 \alpha=\frac{32}{22+5c}.
\end{equation}
Consequently the derivative coefficient is
$\alpha/2=16/(22+5c)$.  The mode commutator has coefficient
$16(m-n)/(22+5c)$, which provides an independent normalization check.
\end{proposition}

\begin{proof}
This is the standard Zamolodchikov $\mathcal W_3$ relation in
Bouwknegt--Schoutens~\cite{BowSch93}.  Passing from the pole-two term
$\alpha\Lambda$ and pole-one term $(\alpha/2)\partial\Lambda$ to modes
produces $(\alpha/2)(m-n)\Lambda_{m+n}$, hence the displayed
commutator coefficient.
\end{proof}

\section{$\mathcal W_3$: exact OPE lines and open scalar projections}
\label{sec:w3-s3-s4-closed-forms}

\begin{theorem}[$T$-line comparison;
\ClaimStatusConditional]
\label{thm:w3-s3-s4-tline}
The $T$-line of $\mathcal W_3$ is the Virasoro OPE at central charge
$c$, with exact norm data
\begin{equation}\label{eq:w3-tline}
 (N_T,N_\Lambda)=\left(\frac c2,\frac{c(5c+22)}{10}\right).
\end{equation}
Under $H_{\mathrm{bar}}+H_{\mathrm{proj}}$, its scalar shadow is the
image of the Virasoro restricted Maurer--Cartan tensor.  The inverse
norm $10/[c(5c+22)]$ is an exact rational function; its identification
with a quartic shadow coefficient is part of the projection theorem.
\end{theorem}

\begin{theorem}[$W$-line comparison;
\ClaimStatusConditional]
\label{thm:w3-w-line-s4-zamolodchikov}
The involution $W\mapsto-W$ preserves the $\mathcal W_3$ OPE, and the
exact $WW\to\Lambda$ coupling is $32/(22+5c)$.  A scalar quartic
coefficient is therefore defined only after the full four-point tensor,
its $TT$, $T\Lambda$, and $\Lambda\Lambda$ channels, and the projection
map have been supplied:
\begin{equation}\label{eq:w3-wline}
 S_4(\mathcal W_3,W)
 :=P_W^{(4)}\bigl(\Theta_{\mathcal W_3}^{(4)}\bigr),
 \qquad
 \text{status: open under }H_{\mathrm{bar}}+H_{\mathrm{proj}}.
\end{equation}
\end{theorem}

\begin{remark}[OPE denominator and shadow denominator]
\label{rem:w3-w-line-denominator-pattern}
The factors $c$, $5c+22$, and $22+5c$ arise respectively from the
$W$ norm, the $\Lambda$ norm, and the $WW\Lambda$ coupling.  A
denominator for $S_4$ also depends on the inverse Gram matrix of every
four-point channel retained by $P_W^{(4)}$.  This gives a concrete
finite computation for the open projection problem.
\end{remark}

\section{Bershadsky--Polyakov: standard conformal data}
\label{sec:bp-s3-s4-closed-forms}

\begin{definition}[Standard BP central charge;
\ClaimStatusDefinitional]
\label{def:bp-arakawa-central-charge}
For the ordinary Bershadsky--Polyakov vertex algebra of
Fehily--Kawasetsu--Ridout~\cite{FKR20}, set
\begin{equation}\label{eq:bp-arakawa-central-charge}
 c_{\mathrm{BP}}(k)
 =-\frac{(2k+3)(3k+1)}{k+3}.
\end{equation}
The four strong generators $J,G^+,G^-,T$ are even, of weights
$1,3/2,3/2,2$.
\end{definition}

\begin{theorem}[BP $T$-line inverse norm;
\ClaimStatusConditional]
\label{thm:bp-t-line-rational-k}
The $T$-line is Virasoro at $c=c_{\mathrm{BP}}(k)$.  Its exact
level-four inverse norm is
\begin{equation}\label{eq:bp-tline-s3-s4}
 \frac{10}{c_{\mathrm{BP}}(k)
 (5c_{\mathrm{BP}}(k)+22)}
 =\frac{10(k+3)^2}
 {3(2k+3)(3k+1)(10k^2+11k-17)}.
\end{equation}
Under $H_{\mathrm{bar}}+H_{\mathrm{proj}}$, this rational function
enters the $T$-line quartic projection.  Cross-channel terms from
$J,G^+,G^-$ remain part of the full BP tensor.
\end{theorem}

\begin{proof}
The identities
\[
 c_{\mathrm{BP}}(k)=-\frac{(2k+3)(3k+1)}{k+3},
 \qquad
 5c_{\mathrm{BP}}(k)+22
 =-\frac{3(10k^2+11k-17)}{k+3}
\]
give the displayed formula by multiplication.
\end{proof}

\begin{theorem}[BP charged-line OPE packet;
\ClaimStatusProvedElsewhere]
\label{thm:bp-other-lines}
In the FKR convention,
\[
 J_{(1)}J=\frac{2k+3}{3},
 \qquad
 G^+_{(2)}G^-=(k+1)(2k+3),
 \qquad
 G^+_{(1)}G^-=3(k+1)J,
\]
and
\[
 G^+_{(0)}G^-
 =3{:JJ:}+\frac{3(k+1)}2\partial J-(k+3)T.
\]
The reverse products follow from ordinary vertex-algebra skew
symmetry because $G^\pm$ are even.  These relations determine the OPE
packet.  Their bar curvature and scalar shadow require
$H_{\mathrm{bar}}+H_{\mathrm{proj}}$.
\end{theorem}

\begin{corollary}[BP reflected central sum;
\ClaimStatusProvedHere]
\label{cor:bp-feigin-frenkel-complementarity}
The critical level reflection $k\mapsto-k-6$ satisfies
\[
 c_{\mathrm{BP}}(k)+c_{\mathrm{BP}}(-k-6)=50.
\]
The shifted secondary expression
$2-24(k+1)^2/(k+3)$ has reflected sum $196$ and belongs to a separate
ledger.  The genus-one quantities
$\kappa_{\mathrm{BP}},\varrho_{\mathrm{BP}},K^\kappa_{\mathrm{BP}}$
remain open pending the complete minimal-DS curvature calculation.
\end{corollary}

\begin{proof}
Substitution into~\eqref{eq:bp-arakawa-central-charge} and cancellation
of the common denominator gives $50$.
\end{proof}

\section{$\mathcal W_\infty$: coordinate and projection data}
\label{sec:w-infinity-psi-degeneration}

\begin{theorem}[$\mathcal W_\infty$ endpoint comparison problem;
\ClaimStatusConditional]
\label{thm:w-infinity-psi-degeneration}
Choose a two-parameter universal $\mathcal W$ algebra, a coordinate
$\Psi$ on its parameter surface, and a normalized weight-three field.
Assume a comparison map to the finite $\mathcal W_N$ truncation loci
and a flat line projection $P_W^{(4)}$.  Then the endpoint quartic is
the specialization
\begin{equation}\label{eq:w-infinity-s4}
 S_4(\mathcal W_\infty[\Psi],W)
 :=P_{W,\Psi}^{(4)}\bigl(\Theta_{\mathcal W_\infty[\Psi]}^{(4)}\bigr).
\end{equation}
Computing this function requires the coordinate conversion from
$\Psi$ to the standard universal OPE parameter and the full
weight-four Gram matrix.  These data constitute the endpoint
comparison package $H_{\mathcal W_\infty}^{(4)}$.
\end{theorem}

\begin{remark}[Finite-rank truncation]
\label{rem:gaberdiel-gopakumar-truncation}
A truncation morphism of universal $\mathcal W$ algebras induces a
line comparison after the weight-three generators and two-point forms
are matched.  This normalization map is the first obligation in
$H_{\mathcal W_\infty}^{(4)}$.
\end{remark}

\section{Super-Yangian $\mathfrak{gl}(1|1)$: a typed parity line}
\label{sec:sl11-fermionic-line-sign}

\begin{theorem}[Parity-line comparison;
\ClaimStatusConditional]
\label{thm:super-yangian-fermionic-line-sign}
Let $V=\mathbb C^{1|1}$ and choose a chiral realization of
$Y_\hbar(\mathfrak{gl}(1|1))$ together with an invariant form.  The
odd basis lines of $V$ contribute Koszul signs to ordered tensor
permutations, while an even bilinear $\psi\psi^*$ carries the product
of two such signs.  A numerical line curvature requires the chosen
invariant form, the current realization, and the residue map; these
form $H_{\mathfrak{gl}(1|1)}^{\mathrm{line}}$.
\end{theorem}

\begin{remark}[Even bilinear line]
\label{rem:fermionic-bilinear-positive-s3}
The bilinear $\psi\psi^*$ has even parity.  Its cubic coefficient is a
three-point contraction in the selected realization and is therefore
computed after the normalization in
$H_{\mathfrak{gl}(1|1)}^{\mathrm{line}}$ has been fixed.
\end{remark}

\section{Exact denominator factors}
\label{sec:denominator-pattern-hypothesis}

\begin{theorem}[OPE denominator factorization;
\ClaimStatusProvedHere]
\label{thm:denominator-factorization-pattern}
The exact denominator factors on the present OPE surface are:
\begin{enumerate}[label=\textup{(\roman*)}]
\item $N_\Lambda=c(5c+22)/10$ for every Virasoro $T$-line;
\item $\alpha=32/(22+5c)$ and $\alpha/2=16/(22+5c)$ for the
 $\mathcal W_3$ $WW$ OPE;
\item $(2k+3)(3k+1)$ and $10k^2+11k-17$ for the BP $T$-line inverse
 norm in~\eqref{eq:bp-tline-s3-s4}.
\end{enumerate}
Further powers in a scalar shadow denominator arise from the inverse
Gram matrix and the contraction graph selected by $P^{(r)}$.
\end{theorem}

\begin{table}[ht]
\centering
\caption{Typed denominator ledger.}
\label{tab:denominator-exponent-pattern}
\renewcommand{\arraystretch}{1.2}
\begin{tabular}{@{}l l l@{}}
\toprule
Line & exact rational datum & scalar shadow status \\
\midrule
Virasoro/$\mathcal W_3$ $T$ & $N_\Lambda=c(5c+22)/10$ & conditional on $P_T^{(4)}$ \\
$\mathcal W_3$ $W$ & $\alpha=32/(22+5c)$ & open four-channel contraction \\
BP $T$ & equation~\eqref{eq:bp-tline-s3-s4} & conditional on BP projection \\
BP $J,G^\pm$ & Theorem~\ref{thm:bp-other-lines} & open mixed-channel tensor \\
$\mathcal W_\infty$ $W$ & universal OPE parameter & open coordinate conversion \\
\bottomrule
\end{tabular}
\end{table}

\section{Extended line census}
\label{sec:extended-census-table}

\begin{table}[ht]
\centering
\caption{Exact line inputs and comparison obligations.}
\label{tab:extended-shadow-census}
\renewcommand{\arraystretch}{1.25}
\begin{tabular}{@{}l l p{5.2cm}@{}}
\toprule
Algebra/line & exact datum & remaining comparison \\
\midrule
$\mathcal W_3/T$ & $N_T=c/2$, $N_\Lambda=c(5c+22)/10$ & scalar quartic projection \\
$\mathcal W_3/W$ & $N_W=c/3$, $\alpha=32/(22+5c)$ & full four-point contraction \\
BP/$T$ & standard $c_{\mathrm{BP}}(k)$ & BP multi-channel projection \\
BP/$J,G^\pm$ & FKR OPE packet & bar curvature and genus-one trace \\
$\mathcal W_\infty/W$ & chosen universal OPE coordinate & endpoint normalization \\
$Y_\hbar(\mathfrak{gl}(1|1))/\psi$ & parity and invariant-form input & chiral realization and residue map \\
\bottomrule
\end{tabular}
\end{table}

\begin{remark}[$E_1$-ordered provenance]
\label{rem:extended-e1-ordered-provenance}
The exact OPE coefficients live at Beilinson level~$1$.  The ordered
bar differential lives at level~$2$ and uses Fulton--MacPherson
residues.  A line shadow is a projection of the level-$2$ or level-$3$
Maurer--Cartan tensor.  The maps in
$H_{\mathrm{line}}+H_{\mathrm{bar}}+H_{\mathrm{proj}}$ connect these
three object types.
\end{remark}

\section{Remaining quartic data}
\label{sec:extended-families-frontier}

\begin{remark}[BP $G^+G^-$ quartic projection; Open]
\label{rem:bp-gpm-closed-form-open}
The FKR products in Theorem~\ref{thm:bp-other-lines} determine the
quadratic input.  The quartic projection requires the $JJ$, $JT$,
$TT$, and charged channels, their Gram matrix, and the signed
configuration-space contractions.
\end{remark}

\begin{remark}[$\mathcal W_\infty$ higher-spin lines; Open]
\label{rem:w-infinity-higher-ws-open}
For $s\ge4$, compute the normalized $W_sW_s$ OPE, match it across the
finite-rank truncation, and construct the projection
$P_{W_s}^{(4)}$.  These three steps determine the first higher-spin
quartic datum.
\end{remark}

\begin{remark}[Super-Yangian $\mathfrak{gl}(m|n)$ lines; Open]
\label{rem:super-yangian-gl-mn-open}
For $m,n\ge0$, the required inputs are a nondegenerate invariant form
on the chosen current presentation, a parity-compatible ordered bar
model, and a scalar trace.  Their construction determines the
bosonic, fermionic, and mixed line curvatures.
\end{remark}
